% ============================================================
% 片段：s03c  对应原文 PDF 第 22--26 页
% 内容：§3.6 镜像几何例（3.6.1 A 模型几何、3.6.2 B 模型几何、3.6.3 求特殊/平坦坐标）
% 译者：t-s03c   审校状态：初稿
% ============================================================

% ==== 原文 PDF 第 22 页 ====

\subsection{镜像几何例}
\label{sec:3.6}

\begin{keytakeaways}
  \begin{itemize}[leftmargin=1.4em,itemsep=2pt,topsep=2pt]
    \item \textbf{物理图像}：五次超曲面 $Z\subset\mathbb{P}^4$ 与其镜像 $Z^*$ 构成镜像对；(3.47) 的 Greene--Plesser 单参数族给出 B 模型一侧的复结构形变族，其模参数 $z$ 与 A 模型的大体积参数 $t$ 由镜像映射 $z(t)$ 相联系。
    \item \textbf{数学工具}：Batyrev 用一对整自反多面体 $(\Delta,\Delta^*)$ 把镜像对实现为环面环绕空间 $(W,W^*)$ 中的超曲面；电荷矩阵 $l_i^a$ 同时编码环面作用、GLSM 的 $U(1)$ 荷与 GKZ 超几何算子。
    \item \textbf{关键结论}：B 模型的周期被 GKZ 扩张系统所湮灭，对五次超曲面约化为四阶 Picard--Fuchs 算子 (3.61)；其判别式 $\Delta=1-3125z$ 的零点给出锥形流形（conifold）点 $z=5^{-5}$，而 $z=0$ 与 $z=\infty$ 分别对应大体积点与轨形点。
    \item \textbf{关键结论}：在 $z=0$ 附近把霍奇滤过的基底取为 (3.67) 并要求联络成为上三角拓扑形式 (3.68)，即可解出镜像映射 $z(t)=q-770q^{2}+\cdots$、规范因子 $S_0(t)$ 以及 $C_{ttt}=\partial_t^3F_0$ 的瞬子展开。
  \end{itemize}
\end{keytakeaways}

以下讨论镜像流形构造的一个实例以及相应镜像映射的求法。这一叙述将相当简略，细节与进一步文献参见 \cite{ref12,ref8,ref48,ref9,ref64}。五次超曲面 $Z$ 指由
\begin{equation}\tag{3.46}
Z:=\{P(x)=0\}\subset\mathbb{P}^4\,,
\end{equation}
定义的 CY 流形，其中 $P$ 是 5 个变量 $x_1,\dots,x_5$ 的 5 次齐次多项式。镜像五次超曲面 $Z^*$ 可用 Greene--Plesser 构造 \cite{ref11} 得到；等价地，也可用 Batyrev 的对偶多面体 \cite{ref65} 在环面几何的语言\footnote{环面几何的综述参见文献 \cite{ref29,ref9}。}中构造。在 Greene--Plesser 构造中，镜像五次超曲面族是由
\begin{equation}\tag{3.47}
\{p(Z^*)=\sum_{i=1}^{5}x_i^5-\psi\prod_{i=1}^{5}x_i=0\}\in\mathbb{P}^4\,,
\end{equation}
给出的单参数五次超曲面族，其中还须取 $(\mathbb{Z}_5)^3$ 商并消解奇点。

下面概述按 Batyrev 的镜像构造。CY 三维流形的镜像对 $(Z,Z^*)$ 由环面环绕空间 $(W,W^*)$ 中的超曲面给出。文献 \cite{ref65} 的镜像对称构造把对 $(Z,Z^*)$ 与一对整自反多面体 $(\Delta,\Delta^*)$ 相联系。

% [Tutor Note: 物理图像] Batyrev 构造中多面体对 $(\Delta,\Delta^*)$ 与镜像对的对应关系
\begin{yizhu}[物理图像]
Batyrev 构造的要点是：$(\Delta,\Delta^*)$ 是互为对偶的整自反多面体，$\Delta$ 的整点给出环面环绕空间 $W$ 的坐标 $u_i$ 与电荷关系 (3.48)，而 $\Delta^*$ 的整点给出镜像 $Z^*$ 的定义方程 (3.51) 中的单项式 $X^{\nu_i}$。“对偶”在物理上就是交换 $h^{1,1}$ 与 $h^{2,1}$：A 模型一侧的凯勒模（环面作用的荷与 GLSM 的 Fayet--Iliopoulos 参数）在 B 模型一侧成为复结构模（$a_i$ 的形变）。因此 (3.52) 中 $n=h^{2,1}(Z^*)=h^{1,1}(Z)$ 并非巧合，而是这一构造的出发点。
\end{yizhu}

% ==== 原文 PDF 第 23 页 ====

\subsubsection{A 模型几何}
\label{sec:3.6.1}

多面体 $\Delta$ 由 $k$ 个相关的整点 $\nu_i$ 刻画，这些整点位于 $\mathbb{Z}^5$ 中一个距原点距离为 $1$ 的超平面内；按文献 \cite{ref65,ref66} 的约定，$\nu_0$ 表示原点。多面体 $\Delta$ 的 $k$ 个整点 $\nu_i(\Delta)$ 对应于环面环绕空间 $W$ 上的齐次坐标 $u_i$，并满足 $n=h^{1,1}(Z)$ 个线性关系：
\begin{equation}\tag{3.48}
\sum_{i=0}^{k-1}l_i^a\nu_i=0\,,\quad a=1,\dots,n\,.
\end{equation}
对固定的 $a$，向量 $l^a$ 的整分量定义了坐标 $u_i$ 在 $\mathbb{C}^*$ 作用下的权 $l_i^a$，即
\[
u_i\to(\lambda_a)^{l_i^a}u_i\,,\quad \lambda_a\in\mathbb{C}^*\,.
\]

$l_i^a$ 也可理解为与上述环面簇 \cite{ref42} 相联系的规范线性 sigma 模型（GLSM）构造中各场的 $U(1)_a$ 荷。环面簇 $W$ 定义为 $W\simeq(\mathbb{C}^k-\Xi)/(\mathbb{C}^*)^n$，其中 $\Xi$ 对应退化轨道的某个例外子集。为构造紧致超曲面，取 $W$ 为某个紧致环面簇上反典范丛的全空间。紧致流形 $Z\subset W$ 通过在 GLSM 中引入超势 $\mathcal{W}_Z=u_0p(u_i)$ 来定义，其中 $x_0$ 是纤维上的坐标，而 $p(u_i)$ 是 $u_{i>0}$ 的多项式，其次数为 $-l_0^a$。
% [原文疑误: 超势写作 $\mathcal{W}_Z=u_0p(u_i)$，纤维坐标却记为 $x_0$，同一坐标出现两种记法；理由：上下文只出现过 $u_0$，且 (3.48) 的坐标均为 $u_i$；建议统一为 $u_0$。此处照抄原文。]
在大凯勒体积下，临界轨迹位于 $u_0=p(u_i)=0$ \cite{ref42}。

CY 流形的一个例子是五次超曲面，即由 $\mathbb{P}^4$ 上反典范丛的一个截面给出的紧致几何。该几何的荷向量为：
\begin{equation}\tag{3.49}
\begin{array}{c@{\quad}c@{\quad}c@{\quad}c@{\quad}c@{\quad}c@{\quad}c@{\quad}c}
 & u_0 & u_1 & u_2 & u_3 & u_4 & u_5 & \\[2pt]
l=\big( & -5 & 1 & 1 & 1 & 1 & 1 & \big)
\end{array}
\end{equation}
多面体 $\Delta$ 的顶点为：
\begin{equation}\tag{3.50}
\begin{array}{l}
\nu_0=(0,0,0,0,0)\,,\quad \nu_1=(1,0,0,0,0)\,,\quad \nu_2=(0,1,0,0,0)\,,\\[2pt]
\nu_3=(0,0,1,0,0)\,,\quad \nu_4=(0,0,0,1,0)\,,\quad \nu_5=(-1,-1,-1,-1,0)\,.
\end{array}
\end{equation}

% [Tutor Note: 易错点] 电荷矩阵 $l_i^a$ 的指标位置与 $u_0/x_0$ 的不一致
\begin{yizhu}[易错点]
(3.48) 中 $l_i^a$ 的上标 $a$ 是荷向量的编号（$a=1,\dots,n$），下标 $i$ 是环面坐标的编号，须与 (3.49) 的表头 $u_0,\dots,u_5$ 逐列对应；五次超曲面有 $n=h^{1,1}=1$，故只有一个荷向量 $l=(-5,1,1,1,1,1)$。核对 (3.48) 时注意求和从 $i=0$ 到 $k-1$，即把原点 $\nu_0$ 也包含在内。此外，正文把超势写成 $\mathcal{W}_Z=u_0p(u_i)$，紧接着却说 $x_0$ 是纤维上的坐标；同一坐标出现 $u_0$ 与 $x_0$ 两种记法（译者按原文照抄），引用 (3.51) 时应按 $u_i$ 统一理解。
\end{yizhu}

\subsubsection{B 模型几何}
\label{sec:3.6.2}

B 模型几何 $Z^*\subset W^*$ 由文献 \cite{ref67,ref65} 的镜像对称构造确定为方程
\begin{equation}\tag{3.51}
p(Z^*)=\sum_{i=0}^{k-1}a_iy_i=\sum_{\nu_i\in\Delta}a_iX^{\nu_i}\,,
\end{equation}
的零点轨迹，其中 $a_i$ 参数化 $Z^*$ 的复结构，$y_i$ 是 $W^*$ 上的齐次坐标 \cite{ref67}，而 $X_m$（$m=1,\dots,4$）是开环面 $(\mathbb{C}^*)^4\subset W^*$ 上的非齐次坐标，且
\[
X^{\nu_i}:=\prod_mX_m^{\nu_i,m}
\]
\cite{ref46}。

% ==== 原文 PDF 第 24 页 ====

关系 (3.48) 对上述齐次坐标施加如下约束：
\begin{equation}\tag{3.52}
\prod_{i=0}^{k-1}y_i^{l_i^a}=1\,,\quad a=1,\dots,n=h^{2,1}(Z^*)=h^{1,1}(Z)\,.
\end{equation}
B 模型中的重要量是全纯 $(3,0)$ 形式，它由下式给出：
\begin{equation}\tag{3.53}
\Omega(a_i)=\mathrm{Res}_{p=0}\,\frac{1}{p(Z^*)}\prod_{i=1}^{4}\frac{\mathrm{d}X_i}{X_i}\,.
\end{equation}
其周期
\begin{equation}\tag{3.54}
\pi_\alpha(a_i)=\int_{\gamma^\alpha}\Omega(a_i)\,,\quad \gamma^\alpha\in H_3(Z^*)\,,\quad \alpha=0,\dots,2h^{2,1}+1\,,
\end{equation}
被 GKZ \cite{ref68} 微分算子的一个扩张系统所湮灭：
\begin{equation}\tag{3.55}
\mathcal{L}(l)=\prod_{l_i>0}\left(\frac{\partial}{\partial a_i}\right)^{l_i}-\prod_{l_i<0}\left(\frac{\partial}{\partial a_i}\right)^{-l_i},
\end{equation}
\begin{equation}\tag{3.56}
\mathcal{Z}_k=\sum_{i=0}^{k-1}\nu_{i,j}\theta_i\,,\quad j=1,\dots,4\,.\qquad
\mathcal{Z}_0=\sum_{i=0}^{k-1}\theta_i+1\,,\quad \theta_i=a_i\frac{\partial}{\partial a_i}\,,
\end{equation}
其中 $l$ 可以是荷向量 $l^a$ 的正整线性组合。方程 $\mathcal{L}(l)\pi_0(a_i)=0$ 由定义 (3.53) 推出。方程 $\mathcal{Z}_k\pi_\alpha(a_i)=0$ 表达周期积分在环面作用下的不变性，并意味着周期积分只依赖于参数 $a_i$ 的特定组合：
\begin{equation}\tag{3.57}
\pi_\alpha(a_i)\sim\pi_\alpha(z_a)\,,\quad z_a=(-)^{l_0^a}\prod_i a_i^{l_i^a}\,,
\end{equation}
其中 $z_a$，$a=1,\dots,n$，定义了 $Z^*$ 的复结构模空间 $\mathcal{M}$ 上的局部坐标。

(3.49) 中定义 A 模型几何的荷向量给出如下镜像几何：
\begin{equation}\tag{3.58}
p(Y)=\sum_{i=0}^{5}a_iy_i=0\,,
\end{equation}
其中坐标 $y_i$ 满足关系
\begin{equation}\tag{3.59}
y_1y_2y_3y_4y_5=y_0^5\,.
\end{equation}
作坐标变换 $y_i=x_i^5$，$i=1,\dots,5$，即显示 (3.47) 与 (3.58) 的等价性，并有
\begin{equation}\tag{3.60}
\psi^{-5}=-\frac{a_1a_2a_3a_4a_5}{a_0^5}=:z\,.
\end{equation}

% ==== 原文 PDF 第 25 页 ====

此外，可以找到如下湮灭 $\tilde\pi_\alpha(z_i)=a_0\pi_\alpha(a_i)$ 的 Picard--Fuchs（PF）算子：
\begin{equation}\tag{3.61}
\mathcal{L}=\theta^4-5z\prod_{i=1}^{4}(5\theta+i)\,,\quad \theta=z\frac{\mathrm{d}}{\mathrm{d}z}\,.
\end{equation}

% [Tutor Note: 易错点] 从 GKZ 扩张系统到四阶 Picard--Fuchs 算子 (3.61)
\begin{yizhu}[易错点]
(3.55) 的 $\mathcal{L}(l)$ 只用到荷向量 $l$ 的正、负分量，是 GKZ 超几何系统的“主”算子；仅它不足以定出周期，还须加上 (3.56) 的 Euler 算子 $\mathcal{Z}_k$（环面作用的不变性）与 $\mathcal{Z}_0$（标度不变性），这正是作者称其为“扩张（extended）系统”的含义，两组方程合起来才把周期约化为单变量 $z_a$ 的超几何函数 (3.57)。对五次超曲面只有一个荷向量 $l=(-5,1,1,1,1,1)$，(3.55) 化为 (3.61) 的 $\theta^4-5z\prod_{i=1}^{4}(5\theta+i)$。可用指标方程核对：$z=0$ 处首项为 $\theta^4$，指标方程 $\rho^4=0$ 给出四重根 $\rho=0$，与 (3.62) 的三个正则奇点 $z=0,\,5^{-5},\,\infty$ 相容，也解释了 $z=0$ 附近的解为何含 $\log z$ 的各次幂；系数 $5$ 与 $\prod(5\theta+i)$ 是超几何级数 $(5n)!/(n!)^5$ 的标准形，改变归一化会同时改变 $z$ 的定义 (3.60) 与判别式中的 $3125=5^5$。
\end{yizhu}

该算子的判别式为
\begin{equation}\tag{3.62}
\Delta=1-3125z\,.
\end{equation}
而汤川耦合可以算出：
\begin{equation}\tag{3.63}
C_{zzz}=\frac{5}{z^3\Delta}\,.
\end{equation}
该 PF 算子给出一个微分方程，它有三个正则奇点，分别对应五次超曲面族的模空间中定义方程变为奇异或获得额外对称性的点，这些点为\footnote{更多细节参见 \cite{ref48}；关于其物理上的 SCFT 解释参见 \cite{ref69}。}：
\begin{itemize}[leftmargin=1.4em,itemsep=1pt,topsep=2pt]
  \item $z=0$，此值处的五次超曲面对应于 $\prod_{i=1}^{5}x_i=0$ 的商流形，它是最退化的 Calabi--Yau 流形，转换到 A 模型一侧对应大半径。
  \item $z=5^{-5}$，这对应微分方程 (3.61) 的一个判别轨迹，也对应定义方程的 Jacobian 消失的轨迹。这类奇点称为锥形流形奇点（conifold singularity）。其 A 模型解释在以后讨论高亏格振幅时将非常有用。
  % [原文疑误: 原文作 "in the context of higher genus amplitudes in at a later stage"，第二个 in 为衍字；译文按 "in the context of higher genus amplitudes at a later stage" 处理。]
  \item $z=\infty$，这称为五次超曲面模空间中的轨形点，对应一个具有大自同构群的非奇异 CY 三维流形；它反映为 5 阶的单值变换（monodromy）。
\end{itemize}

% [Tutor Note: 物理图像] 判别式 $\Delta=1-3125z$ 的零点与三个极限
\begin{yizhu}[物理图像]
$\Delta=1-5^5z$ 是 (3.61) 的判别式。其零点 $z=5^{-5}$（即 $\psi=5$）是锥形流形点：此时定义方程 (3.47) 存在非零解使 $p=\partial p=0$，Jacobian 退化。$z=0$（$\psi\to\infty$）是大复结构极限，即 A 模型的大体积极限，此时基本周期 $S_0\to1$；$z=\infty$（$\psi=0$）是轨形点，其单值变换为 5 阶，反映所取 $(\mathbb{Z}_5)^3$ 商群的残余对称性。镜像对称正是把这三个 B 模型复结构退化点与 A 模型的大体积点、锥形流形点、轨形点对应起来。
\end{yizhu}

\subsubsection{求特殊/平坦坐标}
\label{sec:3.6.3}

B 模型几何的复结构形变通常是定义方程（如 (3.47)）的形变。与所关心的物理生成函数相对应的几何量，是作为这些形变参数的函数计算出来的。若要借助到 A 模型的镜像映射把生成函数物理解释为瞬子求和，就需要找到描述 B 模型一侧形变问题的合适坐标；这就是前面讨论过的特殊/平坦坐标。下面简述如何在镜像五次超曲面模空间中 $z=0$ 附近的区域求出这些坐标。该区域是模空间中与 A 模型的大体积相对应的片，也是可以按量子上同调解释的区域。其思路是：通过对 $\Omega$ 作用导数来张成霍奇滤过，并把 PF 方程表述为形变丛上的联络。
% ==== 原文 PDF 第 26 页 ====
通过改变归一化与坐标，可把该联络化为反映手征环性质的上三角形式。传统上，大体积处的奇点结构也可用来确定镜像映射与汤川耦合，例如 \cite{ref48} 中的综述。不过，对形变丛上的联络作变换，可以在模空间的不同片中找出物理坐标，而这在许多应用中是必需的，例如见 \cite{ref61,ref62,ref28}。

选取如下霍奇滤过截面：
\begin{equation}\tag{3.64}
\vec{\Omega}_z(z)=(\Omega(z),\partial_z\Omega(z),\partial_z^2\Omega(z),\partial_z^3\Omega(z))^t\,,
\end{equation}
便可利用 Picard--Fuchs 方程 (3.61) 读出如下联络矩阵
\begin{equation}\tag{3.65}
\partial_z\vec{\Omega}=
\underbrace{
\begin{pmatrix}
0 & 1 & 0 & 0\\[2pt]
0 & 0 & 1 & 0\\[2pt]
0 & 0 & 0 & 1\\[4pt]
\dfrac{120}{z^3\Delta} & -\dfrac{1-15000z}{z^3\Delta} & -\dfrac{7-45000z}{z^2\Delta} & -\dfrac{6-25000z}{z\Delta}
\end{pmatrix}
}_{M_z}
\vec{\Omega}\,,
\end{equation}
其中 $\Delta=1-3125z$。这可以写成
\begin{equation}\tag{3.66}
\nabla_z\vec{\Omega}_z=(\partial_z-M_z)\vec{\Omega}_z=0\,.
\end{equation}
为找到使联络成为拓扑联络的基与坐标，作如下拟设
\begin{equation}\tag{3.67}
\vec{\Omega}_t(z(t))=\left(
\frac{\Omega(z(t))}{S_0(t)},
\partial_t\frac{\Omega(z(t))}{S_0(t)},
C(z(t))^{-1}\partial_t^2\frac{\Omega(z(t))}{S_0(t)},
\partial_t\left(C(z(t))^{-1}\partial_t^2\frac{\Omega(z(t))}{S_0(t)}\right)
\right)^t,
\end{equation}
该向量按
\[
\vec{\Omega}(z(t))_t=A(t)\vec{\Omega}_z\,.
\]
变换。用
\[
M_t=\left(\partial_tA(t)+\frac{\partial z(t)}{\partial t}A(t)\cdot M_z(z(t))\right)\cdot A^{-1}(t)
\]
变换联络矩阵 $M_z$。要求矩阵 $M_t$ 具有形式
\begin{equation}\tag{3.68}
M_t=
\begin{pmatrix}
0 & 1 & 0 & 0\\[2pt]
0 & 0 & C(z(t)) & 0\\[2pt]
0 & 0 & 0 & 1\\[2pt]
0 & 0 & 0 & 0
\end{pmatrix},
\end{equation}
便得到拟设中各函数的微分方程，其解为：
\[
\begin{aligned}
S_0(t)&=1+120q+21000q^2+14115000q^3+\dots\,,\quad q=\mathrm{e}^{-t}\,,\\
z(t)&=q-770q^2+171525q^3+\dots\,,\\
C_{ttt}&=\partial_t^3F_0=5+2875q+4876875q^2+8564575000q^3+\dots\,.
\end{aligned}
\]

% [Tutor Note: 易错点] 联络矩阵 (3.65) 末行的来源与符号约定
\begin{yizhu}[易错点]
(3.65) 末行可逐项验算：用 $\theta^4=z^4\partial_z^4+6z^3\partial_z^3+7z^2\partial_z^2+z\partial_z$ 把 (3.61) 解出 $\partial_z^4\Omega$，则 $w'$ 的系数为 $11875z^2-z\Delta=z(15000z-1)$、$w''$ 的系数为 $23125z^3-7z^2\Delta=-z^2(7-45000z)$、$w'''$ 的系数为 $6250z^4-6z^3\Delta=-z^3(6-25000z)$、$w$ 的系数为 $120z$，再同除以 $z^4\Delta$ 即得 $120/(z^3\Delta)$、$-(1-15000z)/(z^3\Delta)$、$-(7-45000z)/(z^2\Delta)$、$-(6-25000z)/(z\Delta)$。最常见的错误是漏掉 $\theta^4$ 展开式中的系数 $6$、$7$，或把 $\Delta$ 的符号弄反；可与 (3.63) 的 $C_{zzz}=5/(z^3\Delta)$ 交叉核对。
\end{yizhu}

% [Tutor Note: 物理图像] 为什么大体积展开给出镜像映射
\begin{yizhu}[物理图像]
在 $z=0$ 附近，(3.61) 的四个解中有一个全纯的基本周期 $w_0(z)=1+120z+113400z^2+\cdots$（规范因子），另外三个解含 $\log z$、$\log^2z$、$\log^3z$。取平坦坐标 $t\sim\log z$ 后，$z$ 的逆级数就是镜像映射 $z(t)=q-770q^2+171525q^3+\cdots$（$q=\mathrm{e}^{-t}$），而 $S_0(t)=1+120q+21000q^2+14115000q^3+\cdots$ 正是把 $z(t)$ 代入基本周期 $w_0(z)$ 并整理所得，可逐项验算。$C_{ttt}=\partial_t^3F_0$ 的系数应具有 $\sum_d d^3n_dq^d/(1-q^d)$ 的形式：$2875=n_1$ 是五次超曲面上的直线数，$8\times609250+2875=4876875$（$n_2=609250$ 为二次曲线数），$27\times317206375+2875=8564575000$，这是检验镜像映射是否正确的最直接办法。
\end{yizhu}
